\documentclass[addpoints,12pt,A4]{exam}

\printanswers

\title{Convex Optimization: Problem Set 1}

\input{ExSetup}

\begin{document}
\maketitle

Some problems borrowed from Optimization course of Daniel Dadush. 

\begin{questions}


 \question
Recall the definition of convex hull:
\[ \text{conv}(S) := \bigg\{ \sum_{i=1}^{N} \lambda_{i} \vx_{i} \mid \vx_{i} \in S, \lambda \geq 0, \sum_{i=1}^{N} \lambda_{i} = 1 \bigg\} . \]
 
 \begin{enumerate}
 \item Prove that for finite $S, \text{conv}(S) = \cap_{K \supseteq S} K$ where the intersection is over all closed convex $K$ containing $S$. Therefore $\text{conv}(S)$ is the smallest convex set containing $S$. 

 \item Prove Jensen's inequality: for convex $f : \R^{n} \to \R$ and input $\vx = \sum_{i=1}^{N} \lambda_{i} \vx_{i}$ where $\lambda_{i} \geq 0, \sum_{i=1}^{N} \lambda_{i} = 1$
 \[ f(\vx) \leq \sum_{i=1}^{N} \lambda_{i} f(\vx_{i}) .   \]

 \item Show that if $\CC = \text{conv}\{x_{1}, ..., x_{m}\} \subseteq\R^n$ is a compact convex set and $f:\R^n\to\R$ is a convex continuous function then the {\bf supremum} of $f$ over $\CC$ is attained at one of the vertices $x_{j}$. 
 % \item Show that if $\CC\subseteq\R^n$ is a compact convex set and $f:\R^n\to\R$ is a convex function then the {\bf supremum} of $f$ over $\CC$ is attained at an extreme point of $\CC$.
 \end{enumerate}

  \begin{solution}
  \begin{enumerate}
      \item 
      \item 
      \item 
  \end{enumerate}
  \end{solution}


  \question
  Show the following are equivalent and prove any one of them:
  \begin{itemize}
      \item[(a)] Cauchy-Schwarz inequality for $\R^{n}$
      \[ \langle x, y \rangle = \sum_{i=1}^{n} x_{i} y_{i} \leq \sqrt{ \sum_{i=1}^{n} x_{i}^{2}} \sqrt{\sum_{i=1}^{n} y_{i}^{2}} = \|x\|_{2} \|y\|_{2} ;  \]
      \item[(b)] The triangle inequality $\|x+y\|_{2} \leq \|x\|_{2} + \|y\|_{2}$;
      \item[(c)] Convexity of the Euclidean norm $f(x) := \|x\|_{2}$;
      \item[(d)] Convexity of the unit ball $B_{2}^{n} := \{ x \in \R^{n} \mid \|x\|_{2} \leq 1 \}$.
  \end{itemize}

\begin{solution}
\end{solution}





\question
For convex closed $f : \R^{n} \to \R$ and $x \in \text{dom}(f)$:
\begin{itemize}
    \item Show the sub-level set $L := \{y \mid f(y) \leq f(x)\}$ is a convex set. 
    \item Show for any subgradient $g \in \partial f(x)$, $(g, -1)$ gives a \emph{supporting hyperplane} at $(x,f(x))$ for the epigraph
    \[ \text{epi}(f) := \{ (x, t) \mid f(x) \leq t \} .   \]
    \item Show that $g \in \partial f(x)$ gives a \emph{supporting hyperplane} at $x$ for the sub-level set
\[ L := \{ y \in \R^{n} \mid f(y) \leq f(x) \} .   \]
\end{itemize}

\begin{solution}

\end{solution}


\question
Affine and quadratic functions are the most basic convex functions. We will prove some properties about them:
\begin{itemize}
    \item Show $h : \R^{n} \to \R$ is convex \emph{and} concave (i.e. $h$ is convex and $-h$ is convex) iff $h$ is an affine function, i.e. $h(x) = \langle a, x \rangle + b$ for $a \in \R^{n}, b \in \R$.
    
    \item Let $\tilde{q}(x) := \langle x, Q x \rangle + \langle a, x \rangle + b$ for symmetric matrix $Q \in \R^{n \times n}, a \in \R^{n}, b \in \R$. Show $\tilde{q}$ is convex iff $q(x) := \langle x, Q x \rangle$ is convex. (Hint: use first part)
    
    \item Show $q(x) = \langle x, Q x \rangle$ is convex iff 
    \[ \forall v \in \R^{n}: \quad \langle v, Q v \rangle \geq 0 ,   \]
    and similarly show it is strictly convex iff the inequality is strict for any $v \neq 0$. You can use any of the equivalent characterizations of convexity.
    %the $0$-th order definition of convexity. 
    \\ (These conditions are known as positive-semi-definiteness and positive-definiteness, and are denoted $Q \succeq 0, Q \succ 0$.)

    \item Find the optimizer and optimum value of strictly convex quadratic
    \[ \tilde{q}(x) := \langle x, Q x \rangle + \langle a, x \rangle + b .  \]
\end{itemize}

\begin{solution}

\end{solution}


\question
An Ellipsoid is an affine image of the Euclidean ball
    \[ \mathcal{E} = c + A B_{2}^{n} \quad \text{where} \quad B_{2}^{n} := \{ x \in \R^{n} \mid \|x\|_{2} \leq 1\}  , \]
    for some $c \in \R^{n}$ and $A \in \R^{n \times n}$ invertible.
    \begin{itemize}
    \item Let $q(x) := \langle x, Q x \rangle + \langle d, x \rangle + e$ be \emph{strictly convex}. Show any sub-level set
    \[ L_{t} := \{ x \in \R^{n} \mid q(x) \leq t \}  \]
    is either empty or an Ellipsoid (i.e. find $c,A$ such that $L_{t} = c + A B_{2}$). 

    \item Conversely, given Ellipsoid $\mathcal{E} = c + A B_{2}$ as above, find convex quadratic $q(x) := \langle x, Q x \rangle + \langle d, x \rangle + e$ such that 
    \[ \mathcal{E} = \{ x \in \R^{n} \mid q(x) \leq 1 \} .   \]
    
\end{itemize}
\begin{solution}

\end{solution}


\question
Recall the GLS oracle model for convex sets. For convex set $K \subseteq \R^{n}$, the \emph{polar dual} is 
\[ K^{\circ} := \{w \in \R^{n} \mid \sup_{x \in K} \langle w, x \rangle \leq 1 \} .   \]
\begin{itemize}
    \item Relate VAL and VIOL oracles for $K$ to MEM and SEP oracles for $K^{\circ}$;
    \item Combine this with cutting plane methods to show a reduction from SEParation to OPTimization. In more detail, given access to OPT oracle for $K$, show how to implement the (approximate) SEParation oracle for $K$. 
\end{itemize}
This may seems counter-intuitive, but there are important settings where OPT is easier to implement than SEP (e.g. when $K$ is the spanning tree polytope of a graph). 

\begin{solution}

\end{solution}








\end{questions}
\end{document}
